\documentclass[book.tex]{subfiles}
\begin{document}

\chapter{Runge-Kutta Method}
\label{chap:runge_kutta}
\noindent\rule{11cm}{0.4pt} \\
\textbf{Prerequisites:} \autoref{chap:calculus}, \autoref{chap:ode} \\
\textbf{Difficulty Level:} ** \\
\noindent\rule{11cm}{0.4pt}
\vspace{5mm}

In this chapter we will learn how to understand truncation errors for Runge-Kutta methods. We will then study the representation of second and fourth order Runge-Kutta methods. We will end by solving the Cahn-Hilliard model using Runge-Kutta's fourth-order method.

\section{Runge-Kutta Methods}

In previous chapters, we have discussed Euler's method which is accurate to the first two terms of a Taylor series expansion. By evaluating the slope at multiple places (instead of just the initial value), a better estimate can be obtained. Runge-Kutta methods use this methodology and achieves the accuracy of a Taylor series approach without requiring the calculation of higher derivatives. Many variations exist but all can be cast in the generalized form of
\begin{align}y_{i+1} = y_{i} + \phi h,\end{align}
where $\phi$ is called an increment function, which can be interpreted as a representative slope over the interval. The increment function can be written in general form as

\begin{align}
\phi = a_{1}k_{1} + a_2k_2 + \dots + a_nk_n
\end{align}
where the $a$’s are constants and the $k$’s are
\begin{align}
	\begin{split}
		k_1 &= f(t_i,y_i) \\
		k_2 &= f(t_i + p_1h,y_i + q_{11}k_1h)\\
		k_3 &= f(t_i + p_2h,y_i + q_{21}k_1h+ q_{22}k_2h)\\	
		\vdots&\\
		k_n &= f(t_i + p_{n-1}h,y_i + q_{n-1,1}k_1h + \dots + q_{n-1,n-1}k_{n-1}h)
		\vspace{15pt}
	\end{split}
\end{align}
where the $p$’s, $q$’s are constants and $f$ is the derivative function $dy/dt$. Notice that the $k$’s are recurrence relationships. That is, $k_1$ appears in the equation for $k_2$, which appears in the equation for $k_3$, and so forth.

Various types of Runge-Kutta methods can be devised by employing different numbers
of terms in the increment function as specified by $n$. It can be noted that the first-order RK method with $n$ = 1 is, in fact, Euler’s method. Once $n$ is chosen, values for the $a$’s, $p$’s, and $q$’s are evaluated by setting the above equation equal to terms in a Taylor series expansion. The second-order RK methods (uses an increment
function with two terms) and the fourth-order RK methods (uses an increment
function with four terms) will be discussed herewith. For second-order methods, because terms with $h^3$ and higher are dropped during the derivation, the local truncation error is O($h^3$) and the global error is O($h^2$), Similarly, for fourth-order methods, the global truncation error is O($h^4$).
%\end{document}

\section{Second-Order Runge-Kutta Methods}

The classical second-order Runge-Kutta method is simply obtained from the Taylor series expansion of $y$. It is simplified as

\begin{align}
y_{i+1} &= y_i + f(t_i,y_i)h + \frac{1}{2!}f'(t_i,y_i)h^2 + O(h^3)	\\
y_{i+1} &= y_i + f(t_i,y_i)\frac{h}{2} + \left (\frac{1}{2}f(t_i,y_i) + \frac{1}{2}f'(t_i,y_i)h \right )h + O(h^3)\\
y_{i+1} &= y_i + \left (\frac{1}{2}f(t_i,y_i) + \frac{1}{2}f(t_i+h,y_i + f(t_i,y_i)h)\right )h + O(h^3)
\end{align}

\begin{marginfigure}
	Using Taylor series expansion,
\begin{align}
	\begin{split}
f(t_i+h,y_i + f(t_i,y_i)h))& = f(t_i,y_i) + \frac{\partial f}{\partial t}h \\
& + \frac{\partial f}{\partial y}f(t_i,y_i)h+ O(h^2) \\
& = f(t_i,y_i) + \frac{\partial f}{\partial y}\frac{\partial y}{\partial t}h \\
& = f(t_i,y_i) + f'(t_i,y_i)h 
\end{split}
\end{align}	
\end{marginfigure}
This variant of the second-order Runge-kutta method is known as the \textit{Heun Method}. It can be re-written as

\begin{align}
y_{i+1} = y_{i} + \left (\frac{1}{2}k_1 + \frac{1}{2}k_2\right )h
\end{align}
where
\begin{align}
k_1 = f(t_i,y_i);\hspace{15pt}	k_2 = f(t_i + h,y_i + k_1h)	
\end{align}
However, other variants of these second order RK-methods exist. The general second-order version of the RK equation $y_{i+1} = y_{i} + \phi h$  is given as
\begin{align}
y_{i+1} = y_{i} + (a_1k_1 + a_2k_2)h
\end{align}
where

\begin{align}
k_1 = f(t_i,y_i) ; \hspace{15pt} k_2 = f(t_i + p_1h,y_i + q_{11}k_1h)		
\end{align}

The values of $a$'s, $p$ and $q$ are obtained by setting the above equation equal to a second-order Taylor series. The first three terms of the Taylor series expansion of y is written as  
\begin{align}
y_{i+1} = y_i + f(t_i,y_i)h + \frac{1}{2!}f'(t_i,y_i)h^2 + O(h^3)	
\end{align}

It can be noted that $dy/dt = f(t,y)$. Also,
\begin{align}
f'(t_i,y_i) =\frac{\partial f}{\partial t} + \frac{\partial f}{\partial y}\frac{\partial y}{\partial t}
\end{align}

Substituting this back into the Taylor series expansion, we get
\begin{align}
y_{i+1} = y_i + f(t_i,y_i)h + \frac{1}{2}\left (\frac{\partial f}{\partial t} + \frac{\partial f}{\partial y}\frac{\partial y}{\partial t}\right )_{|_{(t_i,y_i)}}h^2 + O(h^3)
\end{align}

In a similar fashion, $k_2$ can be expanded as
\begin{align}
k_2 = f(t_i + p_1h,y_i + q_{11}k_1h) = f(t_i,y_i) + p_1h\frac{\partial f}{\partial t} +  q_{11}k_1h\frac{\partial f}{\partial y} + O(h^2)	
\end{align}

Now, $y_{i+1} = y_{i} + (a_1k_1 + a_2k_2)h$ is simplified as
\begin{align}
\begin{split}
y_{i+1} & = y_i + \{a_1f(t_i,y_i) +  a_2(f(t_i,y_i) + p_1h\frac{\partial f}{\partial t} +  q_{11}k_1h\frac{\partial f}{\partial y} + O(h^2))\}h \\
 & = y_i + (a_1+a_2)hf(t_i,y_i) + a_2p_1h^2\frac{\partial f}{\partial t} + a_2q_{11}f(t_i,y_i)h^2\frac{\partial f}{\partial y}+ O(h^3)
\end{split}
\end{align}

 On comparing the above equation with (1), three equations can be derived to evaluate the four unknown constants.
\begin{align}
a_1 + a_2 = 1 ; \hspace{10pt}a_2p_1 = 1/2; \hspace{10pt} a_2q_{11} = 1/2
\end{align}

These equations are \textit{underdetermined}, since we have 3 equations with 4 unknowns. Therefore, a value is assumed for one of the unknowns to determine the other three. Suppose that we specify a value for $a_2$, then the equations can be solved as

\begin{align}
		a_1 = 1 - a2; \hspace{10pt} p_1 = q_{11} = \dfrac{1}{2a_2}
\end{align}
 
Because we can choose an infinite number of values for $a_2$, there are an infinite number of second-order RK methods. Every version would yield exactly the same results if the solution to the ODE were quadratic, linear, or a constant. However, they yield different results when (as is typically the case) the solution is more complicated. Two of the most commonly used and preferred versions apart from the Heun's method are discussed here.

\subsection{The Midpoint Method ($a_2 = 1$)}
In this method, $a_2$ is assumed to be 1 and can be solved for $a_1$ = 0 and $p_1$ = $q_{11}$ = 1/2 to obtain
\begin{align}
y_{i+1} = y_{i} + k_2h
\end{align}
where
\begin{align}
k_1 = f(t_i,y_i); \hspace{15pt}k_2 = f(t_i + \frac{1}{2}h,y_i + \frac{1}{2}k_1h)	
\end{align}

\subsection{Ralston's Method ($a_2 = 2/3$)}
Ralston and Rabinowitz (1978) determined that choosing $a_2$ = 2/3 provides a minimum bound on the truncation error for the second-order RK algorithms. In this method, $a_2$ is assumed to be 2/3 and can be solved for $a_1$ = 1/3 and $p_1$ = $q_{11}$ = 3/4 to obtain
\begin{align}
y_{i+1} = y_{i} + (\frac{1}{3}k_1 + \frac{2}{3}k_2)h
\end{align}
where
\begin{align}
k_1 = f(t_i,y_i);\hspace{15pt}	k_2 = f(t_i + \frac{3}{4}h,y_i + \frac{3}{4}k_1h)	
\end{align}

%\end{document}
\section{3-8th Rule Fourth-Order Runge-Kutta Method}

The most popular RK methods are fourth order. As with the second-order approaches, there
are an infinite number of versions. These can be solved for by using a similar approach discussed in the second-order methods. The following variant is one of the most commonly used form. The 3/8th rule fourth-order RK method is similar to Simpson’s 3/8 rule in case of  ODEs that are a function of \textit{t} alone. It is given as

\begin{align}
y_{i+1} = y_{i} + \frac{1}{8}(k_1 + 3k_2 + 3k_3 + k_4)h
\end{align}
where
\begin{align}
\begin{split}
k_1 &= f(t_i,y_i) \\
k_2 &= f(t_i + \frac{1}{3}h,y_i + \frac{1}{3}k_1h)\\
k_3 &= f(t_i + \frac{2}{3}h,y_i - \frac{1}{3}k_1h + k_2h)\\	
k_4 &= f(t_i + h,y_i + k_1h - k_2h + k_3h)
\end{split}
\end{align}

\begin{marginfigure}
	Classical Fourth order Runge-Kutta method:
	\vspace{-10pt}
	\begin{equation*}
	y_{i+1} = y_{i} + \frac{1}{6}(k_1 + 2k_2 + 2k_3 + k_4)h
	\vspace{-6pt}
	\end{equation*}
	where
	\begin{eqnarray*}
		\begin{split}
			k_1 &= f(t_i,y_i) \\
			k_2 &= f(t_i + \frac{1}{2}h,y_i + \frac{1}{2}k_1h)\\
			k_3 &= f(t_i + \frac{1}{2}h,y_i + \frac{1}{2}k_2h)\\	
			k_4 &= f(t_i + h,y_i + k_3h)
		\end{split}
	\end{eqnarray*}
\end{marginfigure}	
It can also be noted that a variant of third-order Runge-Kutta method is also similar to the Simpson's 1/3 rule. In addition, the fourth-order RK method is similar to the Heun approach in that previous estimates of the slope are considered to come up with an improved average slope for the interval. It can also be observed that each of these $k$'s represent a slope and a weighted average of these is considered to arrive at the final improved slope. Another most commonly used variant, is the \textit{classical fourth-order method}. 

\section{Solving Cahn-Hilliard with the Runge-Kutta Method}
Let's consider the Cahn-Hilliard Equation which is an example of a non-linear differential equation describing the process of phase separation, by which the two components of a binary fluid spontaneously separate and form domains which are pure in each component. The Cahn-Hilliard Equation is given as
\begin{align}
\dfrac{\partial c}{\partial t} = D \nabla^2 (c^3 - c - \gamma \nabla^2c )
\end{align} 
Let's solve the Cahn-Hilliard equation given above using the 3/8th rule fourth order Runge-Kutta's method from $t=0$ to $t=5$ with initial conditions as randomized values of -1 and 1 scattered over a grid. We use a time step $(h)$ of $0.002$ s, a diffusion coefficient of $20$ and gamma of $0.5$.

%  d_c = d2x(Conc,dx)+d2y(Conc,dy)
%  d_c3 = d2x(Conc.^3,dx)+d2y(Conc.^3,dy)
%  d_gamma_d_c = d2x(gamma*d_c,dx)+ \
%			d2y(gamma*d_c,dy)
%  dc_dt_out = D*(d_c3 - d_c - d_gamma_d_c)
%  return dc_dt_out
\begin{lstlisting}[language=python]
def $\frac{\partial c}{\partial t}$(c, dx, dy, D, gamma):
  return D*$\nabla^2$(c$^3$-c-$\gamma$*$\nabla^2$c, dx, dy)
\end{lstlisting}
Note that above, by specifying \lstinline{dx,dy}, we are constraining the Laplacian $\nabla^2$ to use a numeric approximation rather than the analytical calculation. In particular, we use a central difference to approximate the second-order derivative at a grid element. We start by initializing a grid.
% x_N = 150, y_N = 150, dx = 1, dy = 1
\begin{lstlisting}[language=python]
def initialize_grid(N: $\mathbb{N}$) $\to \mathbb{R}^{N \times N}$:
  C = zeros(N, N)
  for i j:
    temp = randi(2)
    if (temp==1):
      C[j,i] = -1
    else 
      C[j,i] = 1
C: $\mathbb{R}^{N, N}$ = initialize_grid(N)
\end{lstlisting}

And then we can define a general Runge-Kutta operation
\begin{lstlisting}[language=python]
def runge_kutta($\gamma$ : $\mathbb{R}$, D: $\mathbb{R}$, dt: $\mathbb{R}$, $t_0$: $\mathbb{R}$, $t_e$: $\mathbb{R}$):
  t = $t_0$
  while t < $t_e$:
    k1 = $\frac{\partial c}{\partial t}$(C, dx, dy, D, $\gamma$)  
    $C_2$ =  C + dt/3*k1
    k2 = $\frac{\partial c}{\partial t}$($C_2$, dx, dy, D, $\gamma$)
    $C_3$ =  C + 2*dt/3*k2
    k3 = $\frac{\partial c}{\partial t}$($C_3$, dx, dy, D, $\gamma$)
    $C_4$ =  C + dt*k3
    k4 = $\frac{\partial c}{\partial t}$($C_4$, dx, dy, D, $\gamma$)
    C = C + dt*(k1 + 3*k2 + 3*k3 + k4)/8
    C = bound_conc(C)
    t = t + dt
\end{lstlisting}

The function $\frac{\partial c}{\partial t}$ estimates the derivative at the specified concentration.  The function \lstinline{bound\_conc}() ensures that the concentration lies within [-1,1]. A slightly generalized version of the equation above can be constructed that applies Runge-Kutta to any differential equation system rather than the Cahn-Hilliard equation.
\end{document}
